\chapter{Customer Service Automation with LLMs (Week 11)}

\section{Learning objectives}
\begin{itemize}
  \item Explain how LLMs change the design space of customer service automation in e-commerce.
  \item Distinguish when to use rules, classical NLP, retrieval, and LLM-based systems.
  \item Design a retrieval-augmented generation (RAG) assistant with tool use, workflow orchestration, and escalation paths.
  \item Evaluate LLM service systems using operational, quality, and risk-oriented criteria rather than only surface fluency.
\end{itemize}

\section{Why customer service is a natural LLM use case}
Customer service in e-commerce involves large volumes of language-centered work:
\begin{itemize}
  \item answering policy questions,
  \item helping customers track orders,
  \item explaining returns and refunds,
  \item resolving account issues,
  \item supporting product questions,
  \item handling multilingual or multi-channel interactions.
\end{itemize}

Traditional service automation relied on rules, scripted chat flows, FAQ search, and intent classification.
These tools are still useful, but LLMs change what is possible by supporting more flexible language understanding and generation.

However, flexibility is not enough.
Customer support is operationally sensitive.
A fluent but incorrect answer about order status, refund policy, or account action can create real business harm.
This means LLM systems must be designed as controlled service workflows, not just chat interfaces.

\section{Customer service tasks in e-commerce}
Not all service tasks are equally suited to LLMs.
An architectural view should distinguish several task types.

\subsection{Informational tasks}
Examples include:
\begin{itemize}
  \item store policy questions,
  \item shipping timelines,
  \item payment methods,
  \item warranty coverage,
  \item how to initiate a return.
\end{itemize}

These are often strong candidates for retrieval-augmented LLM assistance because the main challenge is interpreting and presenting known information.

\subsection{Transactional tasks}
Examples include:
\begin{itemize}
  \item checking order status,
  \item canceling an order,
  \item changing a shipping address,
  \item opening a support case,
  \item issuing a return label.
\end{itemize}

These tasks require interaction with operational systems.
The LLM cannot safely invent outcomes.
It must use tools or APIs to read or modify real system state.

\subsection{Judgment-heavy tasks}
Examples include:
\begin{itemize}
  \item refund exceptions,
  \item suspicious behavior review,
  \item policy edge cases,
  \item emotionally complex complaints,
  \item high-value account escalations.
\end{itemize}

These tasks often require human oversight even if an LLM assists with summarization, drafting, or routing.

\section{Rules, classical NLP, and LLMs}
A common mistake is to assume that LLMs should replace all earlier automation methods.
That is rarely the correct design.

\subsection{When rules are best}
Rules remain useful when:
\begin{itemize}
  \item the workflow is deterministic,
  \item the action must be tightly controlled,
  \item compliance requires explicit logic,
  \item the trigger patterns are narrow and stable.
\end{itemize}

Examples include authentication checks, eligibility thresholds, or exact routing conditions.

\subsection{When classical NLP may still be enough}
Classical NLP or simple classifiers can still be appropriate for:
\begin{itemize}
  \item intent routing,
  \item spam detection,
  \item language identification,
  \item sentiment tagging,
  \item FAQ ranking baselines.
\end{itemize}

\subsection{When LLMs are most valuable}
LLMs are strongest when the task requires:
\begin{itemize}
  \item flexible language interpretation,
  \item synthesis from multiple knowledge sources,
  \item paraphrasing and explanation,
  \item conversation continuity,
  \item summarization of long contexts,
  \item structured tool selection under guidance.
\end{itemize}

The right architecture therefore combines all three approaches rather than treating them as mutually exclusive.

\section{Retrieval-augmented generation (RAG)}
RAG is one of the most important patterns for customer service with LLMs.
Instead of asking the model to rely on parametric memory alone, the system retrieves relevant documents or records and asks the model to ground its answer in them.

\subsection{Why RAG matters}
RAG helps because:
\begin{itemize}
  \item support knowledge changes frequently,
  \item policy and catalog information must remain current,
  \item grounded answers are easier to audit,
  \item hallucination risk is reduced when good retrieval is present.
\end{itemize}

\subsection{Knowledge sources}
Typical RAG sources in e-commerce support include:
\begin{itemize}
  \item return and shipping policies,
  \item help-center content,
  \item seller or marketplace rules,
  \item product documentation,
  \item order and shipment records,
  \item CRM case history,
  \item internal support playbooks.
\end{itemize}

\subsection{Retrieval quality}
RAG only works well if retrieval quality is strong.
If the wrong policy document or stale order record is retrieved, the LLM may produce a polished but misleading answer.
This is why retrieval architecture matters as much as model prompting.

\section{Tool use and operational actions}
Customer service systems often need more than retrieved documents.
They need access to operational tools.

\subsection{Examples of service tools}
\begin{itemize}
  \item get order status,
  \item get shipment tracking,
  \item verify return eligibility,
  \item create support ticket,
  \item generate return label,
  \item check account verification state,
  \item retrieve CRM notes.
\end{itemize}

\subsection{Why tools are necessary}
Without tools, the LLM can only talk \emph{about} the workflow.
With tools, it can participate in the workflow while remaining grounded in current system state.

\subsection{Action control}
Not all tools should be available with the same level of autonomy.
Some are read-only and low risk.
Others change customer state and should require:
\begin{itemize}
  \item authentication confirmation,
  \item policy checks,
  \item explicit user confirmation,
  \item or human approval.
\end{itemize}

\section{Workflow orchestration}
A customer service assistant is rarely just a single prompt to a model.
It is more often a workflow with distinct steps.

\subsection{Typical service flow}
A service flow may include:
\begin{enumerate}
  \item classify the issue or infer likely task,
  \item retrieve relevant policy or account context,
  \item call operational tools if needed,
  \item generate a grounded response,
  \item decide whether to resolve, request clarification, or escalate,
  \item log the interaction and outcome.
\end{enumerate}

\subsection{Why orchestration matters}
Without orchestration, an LLM assistant may:
\begin{itemize}
  \item skip required policy checks,
  \item fail to collect missing information,
  \item call tools unnecessarily,
  \item give inconsistent answers across turns,
  \item or fail to escalate when confidence is low.
\end{itemize}

Orchestration is therefore what turns a language model into a service system.

\section{Conversation memory and context}
Customer support conversations often span multiple turns and sometimes multiple channels.
The system must decide what context to preserve.

\subsection{Useful forms of context}
\begin{itemize}
  \item current conversation history,
  \item prior support interactions,
  \item customer account tier,
  \item recent orders and return events,
  \item prior escalations or unresolved issues.
\end{itemize}

\subsection{Context hygiene}
More context is not automatically better.
The system should avoid:
\begin{itemize}
  \item pulling irrelevant records,
  \item mixing incompatible customer identities,
  \item feeding stale or sensitive data without justification,
  \item overloading prompts with low-signal content.
\end{itemize}

Good service design uses enough context to ground the answer without degrading precision or privacy.

\section{Grounding, citations, and answer control}
When the system answers policy or status questions, it should ideally show the basis for the answer.

\subsection{Why citations help}
Citations or source references can:
\begin{itemize}
  \item increase operator trust,
  \item help reviewers audit output,
  \item make it easier to detect retrieval failures,
  \item support handoff to human agents.
\end{itemize}

\subsection{Controlled answer styles}
Service responses often need formatting and policy constraints such as:
\begin{itemize}
  \item concise answer first,
  \item explicit next step,
  \item no unsupported promises,
  \item no invented refund or delivery guarantees,
  \item handoff language when the policy is uncertain.
\end{itemize}

This is one reason prompt design alone is insufficient.
The surrounding system must constrain what the model is allowed to say and do.

\section{Escalation and human handoff}
An LLM service assistant should not aim to resolve every case autonomously.
Escalation is a core capability, not a fallback afterthought.

\subsection{Escalation triggers}
Escalation may be appropriate when:
\begin{itemize}
  \item the policy is ambiguous,
  \item the required action is high risk,
  \item the customer is frustrated or high value,
  \item fraud or compliance concerns are present,
  \item the model lacks sufficient evidence,
  \item repeated attempts have failed.
\end{itemize}

\subsection{Good handoff design}
A good handoff should include:
\begin{itemize}
  \item summary of customer intent,
  \item relevant retrieved documents,
  \item tool outputs already collected,
  \item recommended next action,
  \item uncertainty or policy flags.
\end{itemize}

This can reduce agent workload while preserving human accountability.

\section{Evaluation of LLM service systems}
Evaluation must go beyond subjective fluency.
In customer service, a helpful-sounding answer can still be wrong or risky.

\subsection{Evaluation dimensions}
Useful dimensions include:
\begin{itemize}
  \item factual correctness,
  \item policy compliance,
  \item grounding quality,
  \item task completion success,
  \item escalation appropriateness,
  \item customer satisfaction,
  \item average handling time,
  \item hallucination or unsafe-action rate.
\end{itemize}

\subsection{Offline versus online evaluation}
Offline evaluation may use test conversations, policy QA sets, or synthetic workflows.
Online evaluation may track:
\begin{itemize}
  \item containment rate,
  \item escalation rate,
  \item resolution quality,
  \item recontact rate,
  \item CSAT,
  \item and downstream business costs.
\end{itemize}

\subsection{Human review in evaluation}
Because support quality is nuanced, human review remains important.
Reviewers can assess:
\begin{itemize}
  \item whether the response was grounded,
  \item whether the tone was appropriate,
  \item whether the assistant acted within policy,
  \item whether escalation should have occurred sooner.
\end{itemize}

\section{Guardrails and safety}
Customer service assistants interact with customers directly, so safety controls are essential.

\subsection{Common risk areas}
\begin{itemize}
  \item fabricated order or refund status,
  \item privacy leakage across customer accounts,
  \item unsafe or unauthorized operational actions,
  \item policy violations,
  \item manipulative or inappropriate tone,
  \item prompt injection through retrieved content or user messages.
\end{itemize}

\subsection{Guardrail mechanisms}
Typical guardrails include:
\begin{itemize}
  \item scoped retrieval and identity verification,
  \item tool permissions by action type,
  \item policy filters and post-generation checks,
  \item refusal or defer behavior under uncertainty,
  \item human approval for sensitive workflows,
  \item full logging and audit trails.
\end{itemize}

Guardrails should be treated as first-class design components, not optional moderation add-ons.

\section{Prompt injection and retrieval risks}
When the assistant uses retrieved documents or conversation text, it can be exposed to malicious or misleading content.

\subsection{Examples of injection risk}
\begin{itemize}
  \item a seller document includes hidden instructions not meant for the assistant,
  \item a support message tries to override policy behavior,
  \item external or user-generated content contains adversarial text,
  \item stale or conflicting documents appear in retrieval results.
\end{itemize}

\subsection{Mitigations}
Mitigation strategies include:
\begin{itemize}
  \item restricting retrieval sources,
  \item separating system instructions from retrieved content,
  \item validating tool actions independently of model output,
  \item ranking and curating trusted knowledge bases.
\end{itemize}

\section{A practical architecture for LLM support}
A practical LLM support stack often includes:
\begin{itemize}
  \item a conversation layer,
  \item retrieval over curated support knowledge,
  \item connectors to order, CRM, and return systems,
  \item policy-aware orchestration,
  \item evaluation and monitoring,
  \item escalation into human service tools.
\end{itemize}

\subsection{Read and write separation}
It is often useful to separate:
\begin{itemize}
  \item \textbf{read tools:} retrieve order status, policy, case history,
  \item \textbf{write tools:} cancel order, create refund request, open ticket.
\end{itemize}

Write tools usually deserve stricter safeguards because they change operational state.

\subsection{Role of CRM and ERP integration}
An LLM support system may need:
\begin{itemize}
  \item CRM for case history and customer context,
  \item ERP or commerce systems for order and refund truth,
  \item warehouse or shipping systems for delivery state,
  \item policy repositories for standardized answers.
\end{itemize}

This means the assistant is only as good as its enterprise integration design.

\section{Hands-on lab: a grounded FAQ and service assistant}
The practical goal of this chapter is to move from generic chat to controlled support automation.

\subsection{Lab scenario}
Students are given a small local knowledge base containing return policy, shipping policy, warranty terms, and a mock order-status data source.

\subsection{Lab tasks}
\begin{enumerate}
  \item Build a retrieval-based assistant over the local policy documents.
  \item Add source citations to every policy answer.
  \item Add one read-only tool such as mock order-status lookup.
  \item Define escalation rules for low-confidence or high-risk requests.
  \item Design an evaluation rubric for correctness, grounding, and escalation behavior.
\end{enumerate}

\subsection{Expected learning outcome}
By the end of the lab, students should be able to explain why a useful LLM support assistant is a policy-aware workflow, not just a conversational model.

\section{Managerial implications}
LLM-based service systems can reduce repetitive workload and improve responsiveness, but only if they are embedded in disciplined operational processes.
Organizations should resist the temptation to optimize for containment rate alone.
A service assistant that resolves many tickets badly can be more damaging than one that escalates appropriately.

The correct managerial objective is controlled leverage:
\begin{itemize}
  \item automate routine, well-grounded service work,
  \item support human agents on harder cases,
  \item maintain clear accountability for sensitive actions,
  \item and measure quality continuously.
\end{itemize}

\section{Multiple-choice questions (MCQs)}
\begin{enumerate}
  \item Why is retrieval-augmented generation (RAG) especially valuable for e-commerce customer service?
  \begin{enumerate}
    \item Because it removes the need for any support content.
    \item Because it helps ground answers in current policies, records, and support knowledge rather than relying only on model memory.
    \item Because it guarantees every answer is legally correct.
    \item Because it eliminates escalation.
  \end{enumerate}

  \item Which is the best example of a task that usually requires tool use rather than pure text generation?
  \begin{enumerate}
    \item Summarizing a return policy.
    \item Checking real-time order status for a specific customer.
    \item Rewriting a help-center article title.
    \item Translating a greeting.
  \end{enumerate}

  \item Why should escalation be treated as a core feature in LLM customer service systems?
  \begin{enumerate}
    \item Because the goal is to avoid any automated resolution.
    \item Because some requests are ambiguous, sensitive, or high risk and require human review or approval.
    \item Because escalation improves embedding speed.
    \item Because rules cannot coexist with LLMs.
  \end{enumerate}

  \item Which is the best example of a support-system guardrail?
  \begin{enumerate}
    \item Allowing the model to invent missing refund policies to sound helpful.
    \item Requiring grounded sources, scoped tool permissions, and human approval for sensitive actions.
    \item Removing all logs to protect latency.
    \item Hiding all uncertainty from the customer.
  \end{enumerate}

  \item Why is fluency alone an insufficient evaluation criterion for LLM service assistants?
  \begin{enumerate}
    \item Because customer service is only about grammar.
    \item Because a response can sound natural while still being wrong, ungrounded, or policy-violating.
    \item Because evaluation is unnecessary after deployment.
    \item Because retrieval systems cannot be tested offline.
  \end{enumerate}
\end{enumerate}

\subsection*{Answer key}
\begin{enumerate}
  \item (b)
  \item (b)
  \item (b)
  \item (b)
  \item (b)
\end{enumerate}

\section{Exercises (short)}
\begin{enumerate}
  \item Compare a rule-based FAQ system, a classical intent-routing system, and an LLM+RAG assistant for customer support. State one strength and one limitation of each.
  \item Design a tool schema for a read-only order-status lookup that could be safely used by a support assistant. Include required inputs, expected outputs, and at least two safety checks.
  \item Propose an escalation policy for refund-related requests. Define when the assistant can answer directly, when it can create a request, and when it must hand off to a human agent.
\end{enumerate}

\section{Mini-case (Odoo-linked, vendor-neutral)}
\textbf{Scenario:} A hybrid B2C/B2B company uses a headless storefront, Odoo for orders and invoices, a CRM for support history, and a local knowledge base for policies. The company wants an LLM-based support assistant for chat and agent assistance.

\textbf{Task:} Design the service architecture:
\begin{itemize}
  \item Define which data should be retrieved from policy documents, which should come from tools, and which should remain human-only.
  \item Explain how the assistant would handle order-status questions, return eligibility questions, and refund exceptions.
  \item Identify two guardrails and two evaluation metrics that must be in place before broad deployment.
\end{itemize}
